% preamble.tex - gaya laporan, disalin dari main.tex template proposal
% (Proposal_Color_Quantizer.zip); satu-satunya perubahan: jarak sebelum judul bab (\titlespacing*{\chapter}).
% Tambahan khusus laporan ada di bagian paling bawah.
% Jangan menambah \usepackage: MiKTeX memakai AutoInstall, paket baru bisa terunduh diam-diam.

% ===================== PAKET =====================
\usepackage[a4paper,top=3cm,bottom=3cm,left=3cm,right=2.5cm]{geometry}
\usepackage[utf8]{inputenc}
\usepackage[T1]{fontenc}
\usepackage{mathptmx}
\renewcommand{\ttdefault}{lmtt}
\usepackage{textcomp}
\usepackage{graphicx}
\usepackage[table]{xcolor}
\usepackage{listings}
\usepackage{upquote}
\usepackage{enumitem}
\usepackage{titlesec}
\usepackage{titletoc}
\usepackage{array}
\usepackage{tabularx}
\usepackage{amsmath}
\usepackage{float}
\usepackage[section,above,below]{placeins}
% Gambar boleh bergeser ke halaman berikutnya agar sisa halaman terisi teks
\renewcommand{\topfraction}{0.9}
\renewcommand{\bottomfraction}{0.8}
\renewcommand{\textfraction}{0.07}
\renewcommand{\floatpagefraction}{0.75}
\setcounter{topnumber}{3}
\setcounter{bottomnumber}{2}
\setcounter{totalnumber}{5}
\makeatletter
\setlength{\@fptop}{0pt}
\setlength{\@fpsep}{14pt}
\makeatother
\usepackage{caption}
\usepackage{setspace}
\usepackage{indentfirst}
\usepackage{microtype}
\usepackage{tikz}
\usetikzlibrary{positioning,arrows.meta}
\usepackage[hidelinks]{hyperref}

\graphicspath{{gambar/}}

% ===================== GAYA JUDUL (BAB I, I.1, ...) =====================
\renewcommand{\thechapter}{\Roman{chapter}}
\renewcommand{\thesection}{\thechapter.\arabic{section}}
\renewcommand{\thesubsection}{\thesection.\arabic{subsection}}
\setcounter{secnumdepth}{2}
\setcounter{tocdepth}{1}

\titleformat{\chapter}[display]{\normalfont\normalsize\bfseries\centering}{BAB \thechapter}{10pt}{}
\titleformat{name=\chapter,numberless}[block]{\normalfont\normalsize\bfseries\centering}{}{0pt}{}
% Template proposal memakai 1.2cm (judul bab turun ke sekitar 5,1 cm dari tepi atas);
% laporan memakai 0pt (keputusan Ridlo 28 Nov) agar judul bab sekitar 4 cm dari tepi atas.
\titlespacing*{\chapter}{0pt}{0pt}{8pt}
\titleformat{\section}{\normalfont\normalsize\bfseries}{\thesection.}{0.4em}{}
\titlespacing*{\section}{0pt}{10pt}{2pt}
\titleformat{\subsection}{\normalfont\normalsize\bfseries}{\thesubsection.}{0.4em}{}
\titlespacing*{\subsection}{0pt}{8pt}{2pt}

% ===================== DAFTAR ISI =====================
\renewcommand{\contentsname}{DAFTAR ISI}
\newcommand{\titiktoc}{\titlerule*[0.42em]{.}}
\newcommand{\namabab}{BAB}
\contentsmargin{0pt}
\titlecontents{chapter}[0pt]{}
  {\namabab\ \thecontentslabel\titiktoc\thecontentspage\endgraf\noindent}
  {}
  {\titiktoc\thecontentspage}
\titlecontents{section}[1.2em]{}
  {\thecontentslabel.\enspace}
  {}
  {\titiktoc\thecontentspage}

% ===================== GAMBAR, TABEL, LISTING =====================
\renewcommand{\figurename}{Gambar}
\renewcommand{\tablename}{Tabel}
\renewcommand{\lstlistingname}{Listing}
\captionsetup{font=normalsize,labelfont=normalfont,labelsep=period,justification=centering,skip=6pt}

\definecolor{hijauhdr}{RGB}{226,239,217}
\definecolor{garistabel}{RGB}{166,166,166}
\definecolor{latarkode}{RGB}{245,245,245}
\arrayrulecolor{garistabel}
\newcommand{\hdr}[1]{\cellcolor{hijauhdr}\textbf{#1}}
\newcolumntype{L}[1]{>{\raggedright\arraybackslash}p{#1}}
\newcolumntype{Y}{>{\raggedright\arraybackslash}X}

\newcommand{\kode}[1]{\texttt{\small #1}}

% Penanda bagian yang wajib diisi kelompok (cari "\isi" sebelum submit)
\newcommand{\isi}[1]{\textcolor{red}{[#1]}}

\setlength{\parindent}{1.25cm}
\setlength{\parskip}{6pt}
\setstretch{1.15}
\emergencystretch=3em
\clubpenalty=10000
\widowpenalty=10000
\renewcommand{\arraystretch}{1.35}
\setlist{topsep=2pt,itemsep=2pt,parsep=0pt}

% ===================== TAMBAHAN LAPORAN =====================
% Daftar Gambar dan Daftar Tabel, gaya sama dengan daftar isi.
\renewcommand{\listfigurename}{DAFTAR GAMBAR}
\renewcommand{\listtablename}{DAFTAR TABEL}
\titlecontents{figure}[0pt]{}
  {\figurename~\thecontentslabel.\enspace}
  {}
  {\titiktoc\thecontentspage}
\titlecontents{table}[0pt]{}
  {\tablename~\thecontentslabel.\enspace}
  {}
  {\titiktoc\thecontentspage}

% Nomor persamaan "(II.3)", gambar "III.1", dan tabel "IV.1" mengikuti nomor bab Romawi.
\renewcommand{\theequation}{\thechapter.\arabic{equation}}
\renewcommand{\thefigure}{\thechapter.\arabic{figure}}
\renewcommand{\thetable}{\thechapter.\arabic{table}}

% Warna simpul gambar alur (sama dengan Gambar III.1 proposal): serial, OpenMP, CUDA.
\definecolor{simpulserial}{RGB}{242,242,242}
\definecolor{simpulomp}{RGB}{226,239,217}
\definecolor{simpulcuda}{RGB}{252,228,214}

% Pustaka TikZ tambahan (bagian dari pgf, tidak butuh paket baru) untuk kotak program pada Gambar III.1.
\usetikzlibrary{fit,backgrounds}
